% GazeDiffuse -- EMNLP 2026 Submission
% Template: EMNLP 2023 style (emnlp2023.sty)

\documentclass[11pt]{article}
\usepackage{emnlp2023}
\usepackage{times}
\usepackage{latexsym}
\usepackage{amsmath}
\usepackage{amssymb}
\usepackage{graphicx}
\usepackage{booktabs}
\usepackage{multirow}
\usepackage{tikz}
\usepackage{hyperref}
\usepackage{url}

\title{GazeDiffuse: Controllable Readability in Text Diffusion Models\\via Gaze-Guided Parallel Decoding}

\author{
  Rahil Singhi\textsuperscript{1} \quad
  Siddhant\textsuperscript{1} \quad
  Sai Qian Zhang\textsuperscript{1,2} \\
  \textsuperscript{1}NYU Tandon School of Engineering \quad
  \textsuperscript{2}NYU Courant Institute of Mathematical Sciences \\
  SAI Lab \\
  \texttt{\{rs9174, siddhant, sai.qian.zhang\}@nyu.edu}
}

\begin{document}
\maketitle

% ---------------------------------------------------------------------------
% ABSTRACT
% ---------------------------------------------------------------------------
\begin{abstract}
% TODO: ~150 words. Outline:
% - Readability control in text generation is important for accessibility
%   and education, but prior gaze-guided methods operate autoregressively,
%   applying guidance one token at a time.
% - We propose GazeDiffuse, a training-free method that steers masked
%   diffusion language models (MDLM, LLaDA) at every denoising step using
%   predicted gaze fixation durations from eye-tracking data.
% - Because diffusion models decode all positions in parallel, gaze guidance
%   acts globally, yielding more consistent readability shifts with lower
%   sentence-level variance than autoregressive baselines.
% - Experiments on 50 shared prompts show that GazeDiffuse produces
%   controllable FKGL shifts while preserving fluency (MAUVE) and
%   coherence (self-PPL).
% - GazeDiffuse is the first method combining gaze-guided generation with
%   text diffusion models.
\end{abstract}

% ---------------------------------------------------------------------------
% 1  INTRODUCTION
% ---------------------------------------------------------------------------
\section{Introduction}
\label{sec:intro}

% TODO: Motivate readability control for accessibility and education.
% Introduce the gap: prior gaze-guided work (Sauberli et al., EACL 2026)
% is autoregressive, applying guidance token-by-token.
% State the core idea: applying gaze signals to masked diffusion LMs enables
% parallel guidance across all positions at each denoising step.
% Summarize contributions:
%   (1) GazeDiffuse algorithm,
%   (2) training-free integration with MDLM and LLaDA,
%   (3) empirical evidence of lower FK variance vs. AR baselines.

% ---------------------------------------------------------------------------
% 2  RELATED WORK
% ---------------------------------------------------------------------------
\section{Related Work}
\label{sec:related}

% TODO: Cover three threads and position our contribution at their intersection.

\paragraph{Gaze-Guided Text Generation.}
% TODO: Discuss Sauberli et al. (EACL 2026) and the broader literature on
% using eye-tracking signals for NLP tasks. Contrast AR guidance (one token
% at a time) with our parallel approach.

\paragraph{Discrete Text Diffusion Models.}
% TODO: Review Diffusion-LM (Li et al., 2022), MDLM (Sahoo et al., 2024),
% and LLaDA (Nie et al., 2025). Emphasize masked diffusion as the setting
% for our work and distinguish from continuous diffusion approaches.

\paragraph{Controllable Text Generation.}
% TODO: Briefly survey classifier guidance, PPLM, and other decoding-time
% control methods. Note that most operate in continuous diffusion or AR
% settings; our work is the first discrete-diffusion gaze-guided method.

% ---------------------------------------------------------------------------
% 3  METHOD
% ---------------------------------------------------------------------------
\section{Method}
\label{sec:method}

% TODO: Present the GazeDiffuse algorithm in three parts:
% (1) Background on masked diffusion LMs (forward/reverse process),
% (2) Gaze predictor (BERT fine-tuned on GECO corpus),
% (3) Gaze-guided scoring rule: score = log P_LM(token) + lambda * gaze(token).
% Include the full algorithm box and explain the lambda control parameter.

% Method overview figure placeholder
% TODO: Create a TikZ diagram illustrating the GazeDiffuse pipeline.
%
% Suggested layout (left to right):
%
%   [Prompt + MASK tokens]
%       |
%       v
%   +---------------------------+
%   | Frozen MDLM / LLaDA      |  --> log P_LM(token | context)
%   +---------------------------+
%       |
%       v
%   +---------------------------+
%   | Gaze Predictor (BERT)     |  --> gaze(token) = predicted fixation duration
%   +---------------------------+
%       |
%       v
%   +---------------------------+
%   | Combined Scoring          |
%   | score = log P + lambda *  |
%   |         gaze(token)       |
%   +---------------------------+
%       |
%       v
%   [Reveal highest-confidence tokens]
%       |
%       v  (repeat for T denoising steps)
%   [Final generated text]
%
% Key visual elements:
%   - Show parallel arrows across ALL masked positions (not sequential)
%   - Color-code lambda direction: blue for easier (lambda < 0),
%     red for harder (lambda > 0)
%   - Contrast with a small inset showing AR (left-to-right, one at a time)
%
% Approximate size: \textwidth wide, ~6cm tall

\begin{figure*}[t]
\centering
% TODO: Replace this placeholder box with the actual TikZ diagram.
\fbox{\parbox{0.9\textwidth}{\centering\vspace{4cm}
  \textbf{Method Overview} \\[6pt]
  GazeDiffuse pipeline: Frozen diffusion LM + gaze predictor \\
  with parallel guidance at each denoising step. \\
  See comments in source for diagram specification.
\vspace{4cm}}}
\caption{
  Overview of the GazeDiffuse method. At each denoising step, a frozen
  masked diffusion LM produces token log-probabilities for all masked
  positions simultaneously. A gaze predictor scores each candidate by
  predicted fixation duration. The combined score
  $s = \log P_{\text{LM}}(w) + \lambda \cdot \text{gaze}(w)$ selects
  tokens, with $\lambda < 0$ favoring easier words and $\lambda > 0$
  favoring harder words.
}
\label{fig:method}
\end{figure*}


\subsection{Masked Diffusion Language Models}
\label{sec:method:mdlm}

% TODO: Define the forward masking process and the reverse denoising
% objective following MDLM (Sahoo et al., 2024). Establish notation
% used throughout the paper.

\subsection{Gaze Predictor}
\label{sec:method:gaze}

% TODO: Describe the BERT-base model fine-tuned on the GECO eye-tracking
% corpus (Cop et al., 2017) to predict mean fixation duration per token.
% Mention the 5-token context window and subject-level cross-validation.

\subsection{Gaze-Guided Parallel Decoding}
\label{sec:method:guidance}

% TODO: Present the core GazeDiffuse scoring function. Explain how lambda
% controls readability direction (negative = easier, positive = harder).
% Contrast with AR gaze guidance to highlight the parallel advantage.

% ---------------------------------------------------------------------------
% 4  EXPERIMENTAL SETUP
% ---------------------------------------------------------------------------
\section{Experimental Setup}
\label{sec:experiments}

% TODO: Describe the experimental design covering models, baselines,
% prompts, and evaluation protocol.

\subsection{Models and Baselines}
\label{sec:experiments:models}

% TODO: List all methods compared: Unguided MDLM, AR + Gaze (GPT-2),
% GazeDiffuse (MDLM), GazeDiffuse (LLaDA 8B). Specify model sizes,
% checkpoints, and any frozen/trainable components.

\subsection{Evaluation Metrics}
\label{sec:experiments:metrics}

% TODO: Define each metric and its role:
% FKGL and ARI for readability, MAUVE for distributional similarity,
% self-PPL for coherence, FK sentence variance for consistency.
% Cite Kincaid (1975), Senter & Smith (1967), Pillutla et al. (2021).

\subsection{Data and Prompts}
\label{sec:experiments:data}

% TODO: Describe the 50 shared prompt seeds, the GECO corpus for gaze
% predictor training, and generation hyperparameters (gen_length, steps,
% lambda values in the sweep).

% ---------------------------------------------------------------------------
% 5  RESULTS
% ---------------------------------------------------------------------------
\section{Results}
\label{sec:results}

% TODO: Present main quantitative results. Reference the results table.
% Highlight: (1) FKGL shift with lambda, (2) MAUVE preservation,
% (3) FK variance comparison between GazeDiffuse and AR baseline.

% Main results table: Methods x Metrics
% Matches the structure from README.md

\begin{table*}[t]
\centering
\small
\begin{tabular}{@{}lccccccc@{}}
\toprule
\textbf{Method}
  & \textbf{FKGL} ($\lambda{=}{-}1$)
  & \textbf{FKGL} ($\lambda{=}0$)
  & \textbf{FKGL} ($\lambda{=}{+}1$)
  & \textbf{ARI}
  & \textbf{MAUVE} $\uparrow$
  & \textbf{Self-PPL} $\downarrow$
  & \textbf{FK Var.} $\downarrow$ \\
\midrule
Unguided MDLM
  & ---
  & \textit{TBD}
  & ---
  & \textit{TBD}
  & ---
  & \textit{TBD}
  & \textit{TBD} \\
AR + Gaze (GPT-2)
  & \textit{TBD}
  & \textit{TBD}
  & \textit{TBD}
  & \textit{TBD}
  & \textit{TBD}
  & \textit{TBD}
  & \textit{TBD} \\
\midrule
\textbf{GazeDiffuse (MDLM)}
  & \textit{TBD}
  & \textit{TBD}
  & \textit{TBD}
  & \textit{TBD}
  & \textit{TBD}
  & \textit{TBD}
  & \textit{TBD} \\
\textbf{GazeDiffuse (LLaDA 8B)}
  & \textit{TBD}
  & \textit{TBD}
  & \textit{TBD}
  & \textit{TBD}
  & \textit{TBD}
  & \textit{TBD}
  & \textit{TBD} \\
\bottomrule
\end{tabular}
\caption{
  Main results across readability and quality metrics.
  FKGL columns show Flesch--Kincaid Grade Level at three guidance strengths.
  MAUVE measures distributional similarity to unguided text (higher is better).
  Self-PPL is perplexity under the base LM (lower is better).
  FK Var.\ is sentence-level FKGL variance (lower indicates more consistent readability).
  % TODO: Fill in numbers once experiments complete.
}
\label{tab:results}
\end{table*}


% ---------------------------------------------------------------------------
% 6  ANALYSIS
% ---------------------------------------------------------------------------
\section{Analysis}
\label{sec:analysis}

% TODO: Deeper investigation beyond the main results table.

\paragraph{Readability Consistency.}
% TODO: Analyze FK sentence variance in detail. Show that parallel
% decoding yields more uniform readability across sentences than
% sequential AR guidance.

\paragraph{Lambda Sensitivity.}
% TODO: Plot or table showing FKGL as a function of lambda for each
% method. Discuss the effective range and diminishing returns.

\paragraph{Qualitative Examples.}
% TODO: Include 2-3 cherry-picked generation examples at different
% lambda values showing readable vs. complex output from the same prompt.

% ---------------------------------------------------------------------------
% 7  CONCLUSION
% ---------------------------------------------------------------------------
\section{Conclusion}
\label{sec:conclusion}

% TODO: Summarize the contribution: GazeDiffuse is the first method to
% combine gaze-guided generation with text diffusion models. Recap the
% key empirical findings (controllable FKGL, lower variance, preserved
% fluency). Discuss limitations and future directions (e.g., scaling to
% larger diffusion LMs, multilingual gaze data, user studies).

% ---------------------------------------------------------------------------
% LIMITATIONS
% ---------------------------------------------------------------------------
\section*{Limitations}

% TODO: Discuss limitations honestly: (1) gaze predictor trained on
% English-only GECO data, (2) evaluation on a single prompt set,
% (3) no human evaluation yet, (4) computational overhead of gaze
% scoring at each denoising step.

% ---------------------------------------------------------------------------
% ETHICS STATEMENT
% ---------------------------------------------------------------------------
\section*{Ethics Statement}

% TODO: Address ethical considerations: use of eye-tracking data,
% potential for misuse (deliberately obscuring text), accessibility
% benefits for low-literacy populations.

% ---------------------------------------------------------------------------
% ACKNOWLEDGMENTS
% ---------------------------------------------------------------------------
\section*{Acknowledgments}

% TODO: Thank Prof. Sai Qian Zhang, NYU SAI Lab, NYU HPC team,
% Kuleshov Group (Cornell) for MDLM codebase, Sauberli et al. for
% the gaze-guided generation framework.

\bibliography{references}
\bibliographystyle{acl_natbib}

\end{document}
